At $d=10$, the ordinary split interval covered $91.5\%$ of test responses with a mean width of 5.410. Local weighting raised coverage to $92.7\%$ and width to 5.575. The paired change in width was $+0.165$ (pointwise 95\% $t$ interval $[-0.023,0.353]$). At $d=100$, the corresponding changes were from $93.2\%$ to $94.3\%$ in coverage and from 8.064 to 8.550 in width. The width increased in all 20 paired trials (mean difference $+0.487$, interval $[0.137,0.836]$). The comparisons use the same fitted network and calibration residuals within each trial. The two rules attain different observed coverages, so these results do not measure width efficiency at exactly matched coverage.

In the high-noise subgroup at $d=10$, observed coverage was $85.3\%$ for ordinary split conformal, $87.4\%$ for local weighting and $76.0\%$ for the fitted Gaussian process. At $d=100$, the corresponding estimates were $89.8\%$, $91.3\%$ and $86.0\%$. The weighted interval's mean widths in the low- and high-noise groups were 5.39 and 5.75 at $d=10$, compared with 8.57 and 8.53 at $d=100$. Thus the difference in width across noise groups largely disappeared in the higher-dimensional experiment. The finite-radius cap described in Section~\ref{sec:uq} was used for $1.14\%$ and $1.22\%$ of test predictions, respectively.
